\section{Introduction}\label{sec:introduction}

For a smooth projective complex variety \(X\), an algebraic subvariety
of codimension \(p\) has a rational cohomology class of Hodge type
\((p,p)\).  The rational Hodge conjecture asks whether these classes
span every rational class of that type.  Write
\[
 \Hdg^p(X)=H^{2p}(X,\Q)\cap H^{p,p}(X,\C),
\]
where the intersection is taken inside \(H^{2p}(X,\C)\), and write
\(\CH^p(X)_\Q\) for codimension-\(p\) cycles modulo rational
equivalence with rational coefficients.

\begin{theorem}\label{thm:main}
Let \(m\ge1\), and let \(S_1,\ldots,S_m\) be any projective complex
K3 surfaces.  For \(X=S_1\times\cdots\times S_m\) and every
\(0\le p\le2m\), the cycle-class map
\[
 \cl:\CH^p(X)_\Q\longrightarrow\Hdg^p(X)
\]
is surjective.
\end{theorem}

The factors may repeat, and no relation between their periods or
endomorphism fields is assumed.  The distinction matters: a product
can contain Hodge tensors coupling different transcendental spaces,
while a power of one surface can contain tensors beyond those generated
by pairwise contractions.  The proof treats both phenomena.

For a projective K3 surface \(S\), its transcendental space is
\[
 T(S)=\NS(S)_\Q^\perp\subset H^2(S,\Q),
\]
with the intersection form.  

\subsection{History and the remaining tensor problem}
The Hodge conjecture, formulated in Hodge's 1950 Congress address
\cite{Hodge1952}, asks which rational cohomology classes come from algebraic
cycles.  For a projective K3 surface itself, the Lefschetz $(1,1)$ theorem,
together with the classes of the surface and of a point, answers the question.
Products introduce a different problem: their Hodge classes include tensors
in the transcendental cohomologies of the factors.  Already for two factors,
the nontrivial part concerns Hodge homomorphisms between these cohomologies;
on higher powers, further invariant tensors occur.  We use the customary
background on K3 surfaces, their periods, and their transcendental Hodge
structures as presented in \cite{HuybrechtsK3}.

Kuga and Satake associated an abelian variety to a polarized Hodge structure
of K3 type \cite{KugaSatake}.  Its first cohomology carries a Clifford-algebra
representation, and the original weight-two Hodge structure embeds into a
tensor square of this weight-one structure.  Deligne developed this
construction in families and used it in his proof of the Weil conjecture
for K3 surfaces \cite[Sections~5--6]{DeligneK3}.  Andr\'e subsequently proved
that the Kuga--Satake correspondence is motivated
\cite[Section~7.1]{AndreMotives}.  These results provide the Hodge-theoretic
and motivic relation to abelian varieties; obtaining an algebraic cycle
representing the correspondence is the additional requirement in the
Kuga--Satake Hodge conjecture.  Van Geemen's account explains the Clifford
construction and its relation to exceptional Hodge classes on abelian
varieties \cite{VanGeemenKS}.

A second line of work constructs correspondences between K3 surfaces using
moduli spaces of sheaves.  Mukai's work supplied the fundamental examples
\cite{MukaiBundles}, and Nikulin extended the isometry result using
quadratic-lattice methods \cite{NikulinCorrespondences}.  Buskin proved that every rational Hodge isometry between
the second cohomologies of two projective K3 surfaces is algebraic
\cite[Theorem~1.1]{Buskin}; Huybrechts gave a proof using twisted derived
categories and established an isomorphism of their Chow motives
\cite[Theorems~0.1--0.2]{HuybrechtsIsogenous}.  The isometry hypothesis is
significant: arbitrary Hodge homomorphisms need not preserve the quadratic
forms.  Varesco proved that algebraic Kuga--Satake correspondences also make
Hodge similarities algebraic \cite[Theorem~4.5]{VarescoSimilarities}.

Zarhin determined the Hodge groups of transcendental K3 cohomology in terms
of its endomorphism field, which is either totally real or CM \cite{Zarhin}.
This description explains why a result about squares need not settle the
problem for powers.  In the totally real case, orthogonal invariant theory
produces determinant tensors in addition to the quadratic contractions.
Varesco identified the corresponding relative determinant classes and
proved a precise generation statement
\cite[Theorem~2.16]{VarescoPowers}.  When the endomorphism field is CM, the
invariants are generated in degree two; the work of Ram\'on Mar\'i, together
with algebraicity of Hodge isometries, gives the Hodge conjecture for all
powers \cite{RamonMari,Buskin}; see also
\cite[Theorem~2.2 and Corollary~2.3]{VarescoPowers}.  Here a CM endomorphism
field is the relevant condition; it need not make the whole transcendental
Hodge structure of CM type.

Algebraic Kuga--Satake correspondences have also been constructed for
important special families.  Paranjape treated the K3 surfaces arising
from double covers of the plane branched along six lines \cite{Paranjape},
and Schlickewei used their Kuga--Satake varieties to establish the Hodge
conjecture for their squares \cite[Theorem~2]{Schlickewei}.  Varesco proved an
all-powers theorem for four-dimensional families of generic Picard number
16 with a specified rational transcendental quadratic space, assuming
algebraicity of the Kuga--Satake correspondences throughout the family
\cite[Theorem~0.2]{VarescoPowers}.  The quadratic-space and family
hypotheses in that statement are essential to its scope.

Recent geometric constructions considerably enlarge these cases.
Floccari proved algebraicity of the Kuga--Satake correspondence and the
Hodge conjecture for all powers of the K3 surfaces associated with varieties
of generalized Kummer type \cite[Theorem~1.2]{FloccariKummer}.
Floccari and Fu obtained related results through singular varieties of
O'Grady's six-dimensional deformation type; their construction also gives
products of powers of the abelian, K3, and hyper-K\"ahler varieties
associated with one such singular variety
\cite[Theorem~5.7 and Corollary~5.8]{FloccariFu}.  These theorems illustrate
how geometric relations with abelian varieties control exceptional tensors
on particular families.  For a general list of K3 surfaces, a further
issue is to control the joint Hodge group: algebraicity of the tensors on
each individual factor does not by itself determine the mixed invariants.



\subsection{The proof and its principal constructions}

The Kuga--Satake realization of \(T(S)\) uses the
weight-one Hodge structure on the full even Clifford algebra
\(W_S=C^+(T(S))\), the \(H^1\) of an abelian variety up to isogeny.
For a rational vector \(w\in T(S)\) of nonzero square, the standard
correspondence is the specified Hodge map
\begin{equation}\label{intro:ks-map}
 T(S)(1)\longrightarrow\End(W_S),\qquad
 v\longmapsto\bigl(x\longmapsto vxw\bigr).
\end{equation}
Here \(\Q(1)\) has weight \(-2\).  A polarization converts this into
a degree-two tensor on the square of the abelian variety.
The precise map in \eqref{intro:ks-map} is essential to the argument.

The bridge from geometry to this exact tensor is
Proposition~\ref{prop:realization}.  It starts with a rational
middle-dimensional class killed by the symplectic form on a quartic
times a large standard torus, or on the torus alone.  Under a condition
on its degree-two annihilator, it produces a perfect complex on the
ordinary mirror.
The conclusion controls two different features of the complex:
its Euler pairings with specified test objects and the dimension of
its second self-extension space.  Part~\ref{part:algebraic} first
uses this proposition to prove Theorem~\ref{thm:main};
Part~\ref{part:realization} then proves the proposition.

The force of the realization statement lies in combining its two
conclusions.  In the K3 construction, the Euler pairings initially
determine only the algebraic component of the complex's Mukai vector.
We choose that component
using Clifford multiplication and separated powers of invariant
alternating tensors.  The Hochschild action on the whole vector
factors through \(\Ext^2\), so the small extension space forces
many cohomological operators to annihilate it.  A Clifford commutant
calculation shows that so large a kernel is possible only for the
prescribed transcendental component, with its exact normalization.
Thus an upper bound on self-extensions determines a tensor that the
Euler pairings cannot see.

The same calculation attains equality in the rank bound.  Equality
makes the semiregularity trace injective, so the complex formally
deforms along its Hodge locus.  This uses the deformation
principle developed by
Buchweitz--Flenner and Pridham
\cite{BuchweitzFlennerSemiregularity,PridhamSemiregularity};
Markman's construction of Weil classes by deforming semiregular
sheaves gives a related application
\cite[Sections~7--9]{MarkmanSecant}.  We apply this principle along
the coupled K3 and Kuga--Satake period family.  Algebraization and
specialization give \eqref{intro:ks-map} on every projective K3
surface, as proved in
Theorem~\ref{thm:exact-ks}.

This use of an extension-space bound is useful beyond identifying
one tensor: it obtains the required Hochschild rank on the algebraic
mirror itself.  It requires no identification of the Hochschild
action under the mirror equivalence.  The same method, applied to a
torus alone, constructs the degree-four spin tensors of
Proposition~\ref{prop:auxiliary-spin}.

Two companion results enter the final algebraic argument.  The exact
full-even-Clifford criterion of~\cite[Theorem~1.1]{AKS} gives all
Hodge classes on self-powers once \eqref{intro:ks-map} is algebraic.
The rational Hodge theorem for CM abelian varieties
\cite[Theorem~1.1]{CM} controls the torus constituents of auxiliary
weight-one Hodge structures.  These are theorem inputs with the stated
scope; the proof below supplies the exact KS correspondences and the
additional tensors needed to combine different surfaces.

For the combination, compare the group preserving algebraic
correspondences with the joint Hodge group of the transcendental
spaces.  The self-power results identify the projections of these
groups to each factor; the remaining question is whether the
relations between different factors agree.  Algebraic morphisms
between weight-one abelian realizations
identify the shared simple factors projectively.  Central signs can
still differ, especially among the low-rank orthogonal and unitary
factors sharing \(\mathfrak{sl}_2\).  The degree-four spin tensors
give perfect pairings between local Hom spaces and transport those
signs across a finite grid of representation choices.  A graph
calculation then shows that the remaining signs come from the center
of the joint derived group.  This is the mixed-product step in
Section~\ref{sec:mixed}; self-power algebraicity alone does not
replace it.

It remains to prove the realization proposition.  We represent the
input class, after multiplication, by a spin graded Lagrangian
immersion whose second cohomology is controlled by the class's
degree-two annihilator.  Its self-intersection generators have odd
integer degrees.  A disk-moduli Stokes identity makes the immersed
object unobstructed, and its degree-two endomorphism space is then
bounded by that second cohomology.  Only the embedded test objects
are passed through ordinary mirror symmetry.  A directed comparison
and a categorical projection produce the perfect complex while
retaining the Euler pairings and the extension bound.
The geometry adapts the controlled-immersion construction
of~\cite{A8}.  Sections~\ref{sec:topology}--\ref{sec:realization}
provide the additional arguments for a K3 factor, arbitrary large
torus dimension, integer gradings, the Stokes identity, and the
categorical comparisons.

\subsection{Conventions and organization}

All cycle classes and correspondences have rational coefficients.
An algebraic tensor means a cohomology tensor represented by a rational
algebraic correspondence; no assertion about relations modulo rational
equivalence is implicit.  Hodge groups use the special, or circle,
convention and fix Tate structures.  In calculations with total
cohomology, Tate twists equalizing weights are left implicit only
when the circle action makes their meaning unambiguous.

The symplectic constructions use spin branes, zero relative-spin
background, and integer Maslov gradings.  Novikov fields and their
extensions are embedded in \(\C\) as abstract fields; these embeddings
are not analytic evaluation of a Novikov parameter.  The relevant
mirror objects are fixed before the rational input class is chosen.

In Part~\ref{part:algebraic}, Section~\ref{sec:setup} states the
realization proposition and analyzes its mirror tests.
Section~\ref{sec:cohomology} develops cap-action and deformation tools.
Sections~\ref{sec:exactks} and~\ref{sec:auxiliary} construct algebraic
spin tensors; Section~\ref{sec:mixed} uses them to prove
Theorem~\ref{thm:main}.  Part~\ref{part:realization} proves the proposition:
Sections~\ref{sec:topology} and~\ref{sec:curvature} construct the
immersion and remove its curvature; Sections~\ref{sec:comparison}
and~\ref{sec:realization} compare the categories.
